% GENERATED FILE. Edit canonical lessons, then run npm run generate:books.
% canonical-source-hash: fnv1a64:1e916dbe8637c242
% canonical-lessons: HI-C256-ulti, HI-C256-ghav, HI-C256-patti, HI-C256-tika, HI-C256-goli

\chapter{Vomiting, A wound, A bandage, An injection, A pill}
\label{ch:hi-vomiting-a-wound-a-bandage-an-injection-a-pill}
\hlchaptermodality{256}

\begin{chapteropening}
\textbf{By the end of this chapter:} \emph{I can say vomiting, a wound, a bandage, an injection and a pill.}

Complete the last lesson of chapter 256: I can say vomiting, a wound, a bandage, an injection and a pill.
\end{chapteropening}

\section[\texorpdfstring{\hi{उल्टी} (ulṭī)}{उल्टी (ulṭī)}]{\hi{उल्टी} (ulṭī) — vomiting}

\label{lesson:HI-C256-ulti}



\begin{quote}
Before the new one: say the Hindi for a husband's sister, then the Hindi for a sister's husband.
\end{quote}

\subsection*{You'll want to know: \hi{उल्टी}}
\textbf{\hi{उल्टी}} — \emph{ulṭī} — \textquotedblleft{}vomiting\textquotedblright{}.

\begin{cousinweb}
\textbf{\hi{उल्टी}} (\emph{ulṭī}) is \textquotedblleft{}an overturning\textquotedblright{}, from \emph{ulaṭnā}, \textquotedblleft{}to be overturned\textquotedblright{}, which CDIAL lists under a reconstructed \emph{*ullaṭati}, \textquotedblleft{}turns over\textquotedblright{} (CDIAL 2368).
\end{cousinweb}

\begin{grammarlens}[title={What you've built}]
One more word for this chapter.
\end{grammarlens}

\subsection*{Your turn}
\begin{itemize}
\raggedright
  \item \emph{Say it:} \emph{ulṭī}
  \item \emph{Say it:} \emph{ulṭī}, once more
  \item \emph{Say it:} say \emph{ulṭī}
  \item \emph{Recall:} say the Hindi for a husband's sister, then the Hindi for a sister's husband, then say \emph{ulṭī} again
\end{itemize}

\begin{tcolorbox}[breakable,colback=teal!4,colframe=teal!35!black,title={Before you move on}]
What does \hi{उल्टी} mean? (\textquotedblleft{}Vomiting\textquotedblright{}.) Say it once more.
\end{tcolorbox}
\section[\texorpdfstring{\hi{घाव} (ghāv)}{घाव (ghāv)}]{\hi{घाव} (ghāv) — a wound}

\label{lesson:HI-C256-ghav}



\begin{quote}
Before the new one: say the Hindi for a sister's husband, then the Hindi for vomiting.
\end{quote}

\subsection*{You'll want to know: \hi{घाव}}
\textbf{\hi{घाव}} — \emph{ghāv} — \textquotedblleft{}a wound\textquotedblright{}.

\begin{cousinweb}
\textbf{\hi{घाव}} (\emph{ghāv}) is a tadbhava, a word inherited from Sanskrit and reshaped by sound change. It comes from \emph{ghāta}, \textquotedblleft{}killing; blow, bruise\textquotedblright{} (CDIAL 4460). CDIAL lists the Hindi forms as \emph{ghāu} and \emph{ghāy}.
\end{cousinweb}

\begin{grammarlens}[title={What you've built}]
One more word for this chapter.
\end{grammarlens}

\subsection*{Your turn}
\begin{itemize}
\raggedright
  \item \emph{Say it:} \emph{ghāv}
  \item \emph{Say it:} \emph{ghāv}, once more
  \item \emph{Say it:} say \emph{ghāv}
  \item \emph{Recall:} say the Hindi for a sister's husband, then the Hindi for vomiting, then say \emph{ghāv} again
\end{itemize}

\begin{tcolorbox}[breakable,colback=teal!4,colframe=teal!35!black,title={Before you move on}]
What does \hi{घाव} mean? (\textquotedblleft{}A wound\textquotedblright{}.) Say it once more.
\end{tcolorbox}
\section[\texorpdfstring{\hi{पट्टी} (paṭṭī)}{पट्टी (paṭṭī)}]{\hi{पट्टी} (paṭṭī) — a bandage}

\label{lesson:HI-C256-patti}



\begin{quote}
Before the new one: say the Hindi for vomiting, then the Hindi for a wound.
\end{quote}

\subsection*{You'll want to know: \hi{पट्टी}}
\textbf{\hi{पट्टी}} — \emph{paṭṭī} — \textquotedblleft{}a bandage\textquotedblright{}.

\begin{cousinweb}
\textbf{\hi{पट्टी}} (\emph{paṭṭī}) is a tadbhava, a word inherited from Sanskrit and reshaped by sound change. It comes from \emph{paṭṭa}, \textquotedblleft{}cloth; bandage, fillet\textquotedblright{} (CDIAL 7700). CDIAL thinks the word probably came from a non-Aryan source.
\end{cousinweb}

\begin{grammarlens}[title={What you've built}]
One more word for this chapter.
\end{grammarlens}

\subsection*{Your turn}
\begin{itemize}
\raggedright
  \item \emph{Say it:} \emph{paṭṭī}
  \item \emph{Say it:} \emph{paṭṭī}, once more
  \item \emph{Say it:} say \emph{paṭṭī}
  \item \emph{Recall:} say the Hindi for vomiting, then the Hindi for a wound, then say \emph{paṭṭī} again
\end{itemize}

\begin{tcolorbox}[breakable,colback=teal!4,colframe=teal!35!black,title={Before you move on}]
What does \hi{पट्टी} mean? (\textquotedblleft{}A bandage\textquotedblright{}.) Say it once more.
\end{tcolorbox}
\section[\texorpdfstring{\hi{टीका} (ṭīkā)}{टीका (ṭīkā)}]{\hi{टीका} (ṭīkā) — an injection, a vaccination}

\label{lesson:HI-C256-tika}



\begin{quote}
Before the new one: say the Hindi for a wound, then the Hindi for a bandage.
\end{quote}

\subsection*{You'll want to know: \hi{टीका}}
\textbf{\hi{टीका}} — \emph{ṭīkā} — \textquotedblleft{}an injection, a vaccination\textquotedblright{}.

\begin{cousinweb}
\textbf{\hi{टीका}} (\emph{ṭīkā}) comes from a reconstructed \emph{*ṭikka}, \textquotedblleft{}mark, spot\textquotedblright{}, and CDIAL glosses the Hindi word \textquotedblleft{}caste-mark\textquotedblright{} (CDIAL 5458). Hindi now also uses it for an injection.
\end{cousinweb}

\begin{grammarlens}[title={What you've built}]
One more word for this chapter.
\end{grammarlens}

\subsection*{Your turn}
\begin{itemize}
\raggedright
  \item \emph{Say it:} \emph{ṭīkā}
  \item \emph{Say it:} \emph{ṭīkā}, once more
  \item \emph{Say it:} say \emph{ṭīkā}
  \item \emph{Recall:} say the Hindi for a wound, then the Hindi for a bandage, then say \emph{ṭīkā} again
\end{itemize}

\begin{tcolorbox}[breakable,colback=teal!4,colframe=teal!35!black,title={Before you move on}]
What does \hi{टीका} mean? (\textquotedblleft{}An injection, a vaccination\textquotedblright{}.) Say it once more.
\end{tcolorbox}
\section[\texorpdfstring{\hi{गोली} (golī)}{गोली (golī)}]{\hi{गोली} (golī) — a pill, a bullet}

\label{lesson:HI-C256-goli}



\begin{quote}
Before the new one: say the Hindi for a bandage, then the Hindi for an injection.
\end{quote}

\subsection*{You'll want to know: \hi{गोली}}
\textbf{\hi{गोली}} — \emph{golī} — \textquotedblleft{}a pill, a bullet\textquotedblright{}.

\begin{cousinweb}
\textbf{\hi{गोली}} (\emph{golī}) is a tadbhava, a word inherited from Sanskrit and reshaped by sound change. It comes from \emph{golikā}, \textquotedblleft{}little ball\textquotedblright{} (CDIAL 4321). It is built on \emph{gola}, \textquotedblleft{}ball\textquotedblright{}.
\end{cousinweb}

\begin{grammarlens}[title={What you've built}]
One more word for this chapter.
\end{grammarlens}

\subsection*{Your turn}
\begin{itemize}
\raggedright
  \item \emph{Say it:} \emph{golī}
  \item \emph{Say it:} \emph{golī}, once more
  \item \emph{Say it:} say \emph{golī}
  \item \emph{Recall:} say the Hindi for a bandage, then the Hindi for an injection, then say \emph{golī} again
\end{itemize}

\begin{tcolorbox}[breakable,colback=teal!4,colframe=teal!35!black,title={Before you move on}]
What does \hi{गोली} mean? (\textquotedblleft{}A pill, a bullet\textquotedblright{}.) Say it once more.
\end{tcolorbox}
